\chapter{Lời Thầy Gửi Học Sinh}
\markboth{Lời Thầy Gửi Học Sinh}{A Teacher's Final Words to Students}

\section{Hãy Tin Mình Có Thể Tiến Bộ}
\begin{truyen}{Hãy Tin Mình Có Thể Tiến Bộ}{Believe You Can Improve}
\chuhoa{T}{hầy không cần em phải hoàn hảo ngay. Thầy chỉ mong em đừng vội kết luận rằng mình không thể khá hơn.}
\textit{(I do not need you to be perfect right away. I only hope you do not conclude too early that you cannot become better.)}

Con người tiến bộ nhờ học, sửa và bền bỉ.
\textit{(People improve through learning, correction, and persistence.)}

Khi em tin mình có thể tiến bộ, em sẽ dám làm lại sau khi sai.
\textit{(When you believe you can improve, you dare to try again after mistakes.)}

Niềm tin đúng không phải ảo tưởng. Nó là điểm bắt đầu của nỗ lực thật.
\textit{(A healthy belief is not fantasy. It is the starting point of real effort.)}
\end{truyen}

\begin{baihoc}
Đừng khóa đời mình lại bằng một kết luận vội về giới hạn của bản thân.
\textit{(Do not lock your life with a rushed conclusion about your limits.)}
\end{baihoc}

\begin{ghinhoanh}
\textbf{Short English to remember:}

Improvement begins with believing growth is possible.

\end{ghinhoanh}

\begin{tuhocnhanh}
1. Vì sao cần tin mình có thể tiến bộ? \textit{Why is it important to believe you can improve?}

2. Niềm tin này khác ảo tưởng ở đâu? \textit{How is this belief different from fantasy?}

3. Em đang giới hạn mình bằng kết luận nào? \textit{What conclusion is limiting you right now?}
\end{tuhocnhanh}
\ngancach

\section{Đừng Tự Xem Thường Mình}
\begin{truyen}{Đừng Tự Xem Thường Mình}{Do Not Look Down on Yourself}
\chuhoa{C}{ó những học sinh bị người khác chê ít, nhưng tự chê mình thì rất nhiều.}
\textit{(Some students are criticized little by others, but they criticize themselves a lot.)}

Tự xem thường mình lâu ngày làm em yếu đi từ bên trong.
\textit{(Looking down on yourself for a long time weakens you from within.)}

Thầy muốn em biết rằng em còn đang lớn. Em chưa phải phiên bản cuối cùng của mình.
\textit{(I want you to know that you are still growing. You are not the final version of yourself.)}

Có thể em còn thiếu, nhưng em không hề vô giá trị.
\textit{(You may still lack many things, but you are not without worth.)}
\end{truyen}

\begin{baihoc}
Biết mình chưa giỏi là tỉnh táo. Xem thường mình là bất công với chính mình.
\textit{(Knowing you are not yet strong is clarity. Despising yourself is unfair to yourself.)}
\end{baihoc}

\begin{ghinhoanh}
\textbf{Short English to remember:}

Self-awareness is good; self-contempt is harmful.

\end{ghinhoanh}

\begin{tuhocnhanh}
1. Vì sao không nên tự xem thường mình? \textit{Why should you not look down on yourself?}

2. Khác nhau giữa tự biết mình và tự khinh mình là gì? \textit{What is the difference between self-awareness and self-contempt?}

3. Em cần nói với mình câu nào công bằng hơn? \textit{What fairer sentence do you need to say to yourself?}
\end{tuhocnhanh}
\ngancach

\section{Hãy Sống Trung Thực}
\begin{truyen}{Hãy Sống Trung Thực}{Live Honestly}
\chuhoa{T}{rung thực không làm con đường dễ hơn ngay, nhưng nó làm con người đứng vững hơn.}
\textit{(Honesty does not make the road easier immediately, but it makes a person stand more firmly.)}

Gian dối có thể cho em một cái lợi ngắn, nhưng thường lấy đi sự yên ổn dài.
\textit{(Dishonesty may give a short-term gain, but it often steals long-term peace.)}

Sống trung thực nghĩa là dám nhận điều mình làm, dám sửa điều mình sai, dám học thật chứ không chỉ tỏ ra đang học.
\textit{(Living honestly means owning what you do, correcting what you got wrong, and learning for real instead of only appearing to learn.)}

Thầy mong em giỏi lên trên nền thật, không phải trên một lớp vỏ.
\textit{(I hope you grow strong on something real, not on a shell.)}
\end{truyen}

\begin{baihoc}
Trung thực là phần gốc giữ nhân cách không bị mục từ bên trong.
\textit{(Honesty is the root that keeps character from rotting within.)}
\end{baihoc}

\begin{ghinhoanh}
\textbf{Short English to remember:}

Honesty creates a stable life.

\end{ghinhoanh}

\begin{tuhocnhanh}
1. Vì sao sống trung thực quan trọng? \textit{Why is honest living important?}

2. Cái lợi ngắn của gian dối thường đổi lấy điều gì? \textit{What long-term cost often comes with dishonesty?}

3. Em có đang học thật hay chỉ đang tỏ ra học? \textit{Are you truly learning or only appearing to learn?}
\end{tuhocnhanh}
\ngancach

\section{Hãy Biết Cảm Ơn}
\begin{truyen}{Hãy Biết Cảm Ơn}{Learn to Be Thankful}
\chuhoa{B}{iết ơn không làm em nhỏ đi. Nó làm em sâu hơn.}
\textit{(Gratitude does not make you smaller. It makes you deeper.)}

Khi biết cảm ơn cha mẹ, thầy cô, bạn bè và những người âm thầm giúp mình, em sẽ sống bớt vô tâm hơn.
\textit{(When you thank parents, teachers, friends, and those who quietly help you, you live less carelessly.)}

Người biết ơn thường cũng biết giữ gìn điều tốt hơn.
\textit{(Grateful people also tend to protect good things better.)}

Thầy mong em đừng xem những điều đang có là hiển nhiên mãi mãi.
\textit{(I hope you do not treat what you have as permanent or automatic.)}
\end{truyen}

\begin{baihoc}
Biết ơn giúp con người vừa khiêm tốn hơn, vừa ấm hơn.
\textit{(Gratitude helps people become both humbler and warmer.)}
\end{baihoc}

\begin{ghinhoanh}
\textbf{Short English to remember:}

Gratitude deepens a person.

\end{ghinhoanh}

\begin{tuhocnhanh}
1. Vì sao cần biết cảm ơn? \textit{Why is gratitude important?}

2. Biết ơn làm con người thay đổi thế nào? \textit{How does gratitude change a person?}

3. Hôm nay em cần cảm ơn ai? \textit{Who do you need to thank today?}
\end{tuhocnhanh}
\ngancach

\section{Đừng Sợ Làm Lại}
\begin{truyen}{Đừng Sợ Làm Lại}{Do Not Fear Starting Again}
\chuhoa{C}{ó những lần đời không cho em đi tiếp bằng cách cũ.}
\textit{(There are times when life does not let you continue in the old way.)}

Khi đó, làm lại không phải là nhục.
\textit{(At that point, starting again is not shameful.)}

Nó là biểu hiện của việc em còn muốn cứu điều đúng, còn muốn xây lại phần tốt.
\textit{(It shows that you still want to save what is right and rebuild what is good.)}

Nhiều người mắc kẹt không phải vì họ yếu, mà vì họ quá xấu hổ để bắt đầu lại.
\textit{(Many people remain stuck not because they are weak, but because they are too ashamed to begin again.)}
\end{truyen}

\begin{baihoc}
Làm lại đôi khi không lùi. Nó là bước can đảm nhất để đi tiếp đúng hơn.
\textit{(Starting again is not always moving backward. It can be the bravest way to move forward more truly.)}
\end{baihoc}

\begin{ghinhoanh}
\textbf{Short English to remember:}

Beginning again can be an act of courage.

\end{ghinhoanh}

\begin{tuhocnhanh}
1. Vì sao không nên sợ làm lại? \textit{Why should you not fear starting again?}

2. Điều gì khiến nhiều người ngại bắt đầu lại? \textit{What keeps many people from starting again?}

3. Có việc gì em cần can đảm làm lại không? \textit{Is there anything you need courage to restart?}
\end{tuhocnhanh}
\ngancach

\section{Hãy Giữ Lòng Tử Tế}
\begin{truyen}{Hãy Giữ Lòng Tử Tế}{Keep a Kind Heart}
\chuhoa{T}{hầy không mong em lớn lên thành người chỉ biết hơn thua.}
\textit{(I do not hope you grow into someone who only knows competition.)}

Thầy mong em giữ được lòng tử tế ngay cả khi đời làm mình mệt.
\textit{(I hope you keep kindness even when life wears you down.)}

Người tử tế không phải lúc nào cũng dễ sống, nhưng họ làm đời đáng sống hơn cho người quanh mình.
\textit{(Kind people do not always have an easy life, but they make life more livable for others.)}

Nếu sau này em có giỏi, có thành công, xin đừng đánh rơi phần người ấy.
\textit{(If one day you become capable and successful, please do not lose that human part.)}
\end{truyen}

\begin{baihoc}
Giữ lòng tử tế là giữ cho thành công của mình còn đẹp.
\textit{(Keeping kindness preserves the beauty of success.)}
\end{baihoc}

\begin{ghinhoanh}
\textbf{Short English to remember:}

Kindness keeps success humane.

\end{ghinhoanh}

\begin{tuhocnhanh}
1. Vì sao cần giữ lòng tử tế khi lớn lên? \textit{Why should you keep kindness as you grow?}

2. Thành công thiếu tử tế sẽ thế nào? \textit{What happens when success lacks kindness?}

3. Em muốn giữ phần người nào của mình lâu nhất? \textit{What human quality do you most want to keep?}
\end{tuhocnhanh}
\ngancach

\section{Hãy Quý Gia Đình Và Người Thương Em}
\begin{truyen}{Hãy Quý Gia Đình Và Người Thương Em}{Value Family and Those Who Love You}
\chuhoa{K}{hi còn nhỏ, nhiều em nghĩ gia đình lúc nào cũng ở đó nên dễ quên quý.}
\textit{(When young, many students assume family is always there and forget to value it.)}

Chỉ đến khi xa nhà, va vấp hoặc mệt mỏi, ta mới thấy những người thương mình âm thầm quan trọng đến mức nào.
\textit{(Only when away from home, struggling, or exhausted do we see how important loving people have quietly been.)}

Thầy mong em đừng đợi quá muộn mới dịu với cha mẹ, ông bà hay người thân.
\textit{(I hope you do not wait too long before becoming gentler with parents, grandparents, or loved ones.)}

Biết quý người thương mình là một phần của trưởng thành.
\textit{(Valuing those who love you is part of maturity.)}
\end{truyen}

\begin{baihoc}
Có những điều quý vì ở rất gần, nên con người lại dễ coi nhẹ nhất.
\textit{(Some things are precious precisely because they are near, and that is why people often take them lightly.)}
\end{baihoc}

\begin{ghinhoanh}
\textbf{Short English to remember:}

Do not take loving people for granted.

\end{ghinhoanh}

\begin{tuhocnhanh}
1. Vì sao con người dễ coi nhẹ gia đình? \textit{Why do people often take family lightly?}

2. Khi nào ta thường nhận ra người thương mình quan trọng? \textit{When do we usually realize how important loving people are?}

3. Em có thể quan tâm gia đình hơn bằng việc nhỏ nào? \textit{What small act can show more care to your family?}
\end{tuhocnhanh}
\ngancach

\section{Học Không Chỉ Vì Điểm Số}
\begin{truyen}{Học Không Chỉ Vì Điểm Số}{Learning Is More Than Grades}
\chuhoa{Đ}{iểm số có giá trị của nó, nhưng nó không phải toàn bộ ý nghĩa của việc học.}
\textit{(Grades have value, but they are not the full meaning of learning.)}

Nếu em chỉ học vì điểm, em sẽ rất dễ mệt, rất dễ gian dối, và rất dễ quên mất điều cần giữ lâu.
\textit{(If you study only for grades, you may become tired, dishonest, and forget what matters long term.)}

Học còn để hiểu, để lớn lên, để có nền cho cuộc đời sau này.
\textit{(Learning is also for understanding, growing, and building a foundation for life.)}

Thầy mong em vẫn cố gắng điểm số, nhưng đừng để nó lấy mất linh hồn của việc học.
\textit{(I hope you still work for grades, but do not let them take away the soul of learning.)}
\end{truyen}

\begin{baihoc}
Điểm số nên là kết quả phụ của việc học thật, không phải mục tiêu duy nhất.
\textit{(Grades should be a byproduct of real learning, not the only goal.)}
\end{baihoc}

\begin{ghinhoanh}
\textbf{Short English to remember:}

Grades matter, but learning matters more.

\end{ghinhoanh}

\begin{tuhocnhanh}
1. Vì sao học không nên chỉ vì điểm? \textit{Why should learning not be only for grades?}

2. Nếu chỉ chạy theo điểm số, mình dễ gặp điều gì? \textit{What risks come from chasing grades alone?}

3. Em muốn giữ ý nghĩa nào đẹp hơn của việc học? \textit{What deeper meaning of learning do you want to keep?}
\end{tuhocnhanh}
\ngancach

\section{Thầy Mong Em Sống Thành Người Đàng Hoàng}
\begin{truyen}{Thầy Mong Em Sống Thành Người Đàng Hoàng}{I Hope You Grow Into a Decent Person}
\chuhoa{S}{au cùng, điều thầy mong không chỉ là em đỗ đạt hay thành công.}
\textit{(In the end, what I hope is not only that you pass exams or become successful.)}

Thầy mong em trở thành người đàng hoàng: biết đúng sai, biết thương người, biết làm việc tử tế, biết đứng vững bằng nhân cách của mình.
\textit{(I hope you become a decent person: one who knows right and wrong, cares for others, works properly, and stands firmly on character.)}

Nếu em làm được điều đó, dù đi đường nào, cuộc đời em vẫn có giá trị sâu xa.
\textit{(If you can do that, then no matter which road you take, your life will hold deep value.)}

Còn với thầy, đó đã là niềm vui rất lớn của một người dạy học.
\textit{(And for a teacher, that is already a great joy.)}
\end{truyen}

\begin{baihoc}
Đích đẹp nhất của giáo dục không chỉ là người thành công, mà là người đàng hoàng.
\textit{(The most beautiful goal of education is not only a successful person, but a decent one.)}
\end{baihoc}

\begin{ghinhoanh}
\textbf{Short English to remember:}

Education should help form a decent human being.

\end{ghinhoanh}

\begin{tuhocnhanh}
1. Điều thầy mong cuối cùng ở học sinh là gì? \textit{What is the teacher's deepest final hope for students?}

2. Người đàng hoàng có những nét nào? \textit{What qualities define a decent person?}

3. Em muốn mình trở thành người thế nào sau này? \textit{What kind of person do you want to become?}
\end{tuhocnhanh}
\ngancach

\section{Thầy Tin Em Sẽ Lớn Lên Tốt Đẹp}
\begin{truyen}{Thầy Tin Em Sẽ Lớn Lên Tốt Đẹp}{I Believe You Will Grow Well}
\chuhoa{T}{rước khi khép lại quyển sách này, thầy muốn nói một điều rất giản dị: thầy tin em.}
\textit{(Before closing this book, I want to say something simple: I believe in you.)}

Không phải vì em đã làm mọi thứ thật tốt, mà vì thầy biết con người có thể trưởng thành rất nhiều nếu không bỏ mình.
\textit{(Not because you have already done everything well, but because people can grow greatly if they do not give up on themselves.)}

Sẽ còn lúc em mệt, sai, chậm, hoặc hoang mang.
\textit{(There will still be times when you are tired, mistaken, slow, or confused.)}

Nhưng nếu em giữ được sự học hỏi, tử tế và bền lòng, em sẽ lớn lên theo một cách đáng quý.
\textit{(But if you keep learning, kindness, and steadiness, you will grow in a deeply worthy way.)}
\end{truyen}

\begin{baihoc}
Một lời tin tưởng đúng lúc đôi khi đủ để một người trẻ đi tiếp đoạn đường dài hơn.
\textit{(A timely word of trust can sometimes help a young person continue much farther.)}
\end{baihoc}

\begin{ghinhoanh}
\textbf{Short English to remember:}

Keep learning, keep kindness, and keep going.

\end{ghinhoanh}

\begin{tuhocnhanh}
1. Vì sao một lời tin tưởng có thể quan trọng? \textit{Why can a word of trust be important?}

2. Em muốn giữ lại điều gì nhất sau quyển sách này? \textit{What do you most want to keep from this book?}

3. Câu nào em muốn tự nhắc mình khi thấy nản? \textit{What sentence do you want to tell yourself when discouraged?}
\end{tuhocnhanh}
\ngancach
